\section{Introduction}
Transformer-based autoregressive models require storing Key ($K$) and Value ($V$) tensors across extended sequence lengths to avoid recomputing self-attention matrices during generation. However, standard KV-cache mechanisms scale linearly in memory footprint $\mathcal{O}(L \cdot D)$ with sequence length $L$ and hidden dimension $D$. Under long-context workloads ($L \ge 128\text{k}$), memory consumption quickly exhausts physical GPU vRAM, necessitating heavy quantization or lossy context dropping. Furthermore, storing raw unencrypted activation vectors in physical GPU memory exposes deep LLM deployments to memory inspection attacks and unauthorized state extraction.

\textbf{ResonanceShardKVCache} resolves the dual challenge of memory explosion and security by uniting $\phi^\infty$ Shard Folding compression with the \textbf{Quantum Shadow Veil (QSV)}.

\subsection{Key Architectural Innovations}
\begin{enumerate}[label=\textbf{\arabic*.}]
    \item \textbf{Dynamic Shard Capacity Thresholding}: Active key-value tokens are accumulated up to $N = 1024$ tokens before executing automatic phase-folding compression.
    \item \textbf{Quantum Shadow Veil (QSV) Encryption}: Sharded historical memory blocks are encrypted via Fibonacci pseudo-random keys and golden-ratio phase shifts $(x + \text{salt}) \cdot \phi^{-n}$.
    \item \textbf{Veil-Authenticated Retrieval}: Decryption operations preserve gradient lattice structures while guaranteeing bitwise recovery of folded attention states.
\end{enumerate}
